\setupcolors[state=start]
\setupcolor[hex]
\definecolor[darkblue][h=003366]
\definecolor[darkgray][h=e6e6e6]
\definecolor[redline][h=d9534f]
\definecolor[blueline][h=5bc0de]
\definecolor[greenline][h=5cb85c]

\setuphead[section][style=\bfa, color=darkblue, before=\blank, after=\blank]
\setuphead[subsection][style=\bf, color=black, before=\blank, after=\blank]
\setuppapersize[A4][A4]
\setupbodyfont[11pt]
\setuplayout[width=middle, height=middle, topspace=2cm, bottomspace=2cm]
\mainlanguage[es]

\starttext
\title{\bf Exploración de OPP\_v5\_1\_mono: dónde deja de cerrar el modelo monolítico}
\blank[big]

Este documento es el par de la exploración de OPP\_v5\_2\_split
(\type{docs/thesis/opp_v5_2_split_exploration.tex}). Las preguntas son las mismas, pero
aplicadas al modelo monolítico: dónde deja de probar el óptimo, qué lo encarece y cuánto más
lento es que el split sobre la misma instancia. Reúne los estudios 08 a 11 y cierra con la suite
de stress de los dos modelos.

\blank

\startsection[title={Alcance y fuentes}]
\startitemize
\item {\bf Estudio 08, frontera:} una grilla de lotes (1, 2, 3, 4 y 6) por lado de la grilla
del depósito (3×3 a 7×7), con el perfil realista de los tests de stress: 26 órdenes por lote,
pesos 4/5, dos bandas de velocidad (0,8 y 1,2) y dos lotes por picker. Tope de 10 minutos por
modelo. Cuando el monolítico no prueba el óptimo en un lado, ese lado y los mayores se saltean
para los lotes siguientes.
\item {\bf Estudio 09, casos compartidos:} 15 instancias idénticas a las de los estudios 01,
02 y 03 del split, con el mismo límite de SCIP (300 s). Todas tienen 6 lotes, en 7×7 salvo
una en 10×10.
\item {\bf Estudio 10, factores:} las preguntas de los estudios 01, 02, 03 y 07 del split a una
escala que el monolítico resuelve (4 lotes, 2 pickers, 5×5), moviendo un factor por vez.
\item {\bf Estudio 11, 40 minutos:} los tres puntos donde el monolítico dejó de cerrar en el
estudio 08, más la instancia exacta del test de stress 04, con el mismo tope que la suite.
\item {\bf Puntos previos:} T2 (estudio 04, 5×5 con 4 lotes) y M0 (estudio 05, la instancia
exacta del test de stress 04).
\item {\bf Todos los puntos son pares:} la misma instancia se resuelve primero con el split y
después con el monolítico. Los dos modelos comparten el óptimo, así que el del split es la
referencia contra la que se mide al monolítico: cuándo encontró una solución cercana, cuándo
llegó al óptimo y cuánto le faltaba a su cota dual.
\item {\bf Datos:} \type{experiments/opp_v5/studies/08_mono_frontier/},
\type{09_shared_with_split/}, \type{10_mono_scale_axes/} y \type{11_mono_long_cap/}.
\stopitemize
\stopsection

\startsection[title={Cómo leer los números}]
\startitemize
\item {\bf La máquina estaba ocupada.} Los estudios 08 y 09 corrieron mientras la suite de
stress del split ocupaba unos dos núcleos. Los tiempos absolutos están inflados. Dentro de cada
par, split y monolítico corrieron uno después del otro con la misma carga de fondo, así que la
proporción entre los dos es comparable.
\item {\bf Con un tope fijo, la carga también afecta el resultado del monolítico:} procesa menos
nodos en el mismo tiempo, así que su brecha al corte sale algo peor que con la máquina sola.
\item {\bf Los nodos siguen siendo deterministas} y son la medida justa para comparar
dificultad (ver el documento del split).
\stopitemize
\stopsection

\startsection[title={El modelo monolítico}]
El monolítico resuelve en un solo MIP lo que el split resuelve en dos etapas: las rutas de todos
los lotes, la asignación a pickers y la secuencia, con la misma función objetivo. La diferencia
estructural está en el tiempo de ruta de cada lote, $R_b$. En el split es un dato: la Etapa 1
lo fija en el mínimo antes de programar. En el monolítico es una variable, definida por los
arcos elegidos, que entra en las restricciones big-M de duración y de no superposición.

La equivalencia de los dos modelos se apoya en que la ruta más corta domina a cualquier otra.
El split usa ese hecho de entrada; el monolítico tiene que descubrirlo durante la búsqueda, y
puede explorar combinaciones de rutas que nunca pueden mejorar el objetivo.

El split no usa ningún ajuste propio: los parámetros de SCIP, la función objetivo y las
constantes big-M ($M$ y $M_{seq}$) son las mismas en los dos modelos. Lo que cambia son dos
simplificaciones que la división habilita, porque con la ruta ya conocida el tiempo de ejecución
de cada lote en cada picker es un número:

\blank[medium]
\bTABLE[split=yes, align=middle, frame=on, width=fit]
\setupTABLE[row][1][background=color, backgroundcolor=darkgray, style=bold]
\bTR \bTH \eTH \bTH Split \eTH \bTH Monolítico \eTH \eTR
\bTR \bTD Ruteo \eTD \bTD Etapa 1, un problema por lote \eTD \bTD dentro del mismo MIP \eTD \eTR
\bTR \bTD Tiempo de ejecución ($t_{exec}$) \eTD \bTD dato \eTD \bTD variable, con dos restricciones big-M por par picker-lote \eTD \eTR
\bTR \bTD Cadena de frío \eTD \bTD los pares no elegibles no existen \eTD \bTD cota sobre una variable \eTD \eTR
\bTR \bTD Paralelismo \eTD \bTD opcional, con más de 4 lotes \eTD \bTD -- \eTD \eTR
\eTABLE
\blank[medium]

Las dos filas del medio no son agregados: son consecuencia directa de dividir. El paralelismo
no interviene en casi nada de lo medido acá: los puntos de 1 a 4 lotes no pasan el umbral, y con
6 lotes en 7×7 la Etapa 1 dura segundos y el tiempo se va en la Etapa 2, que no se paraleliza.

\blank[medium]
\bTABLE[split=yes, align=middle, frame=on, width=fit]
\setupTABLE[row][1][background=color, backgroundcolor=darkgray, style=bold]
\bTR \bTH Instancia \eTH \bTH Split, Etapa 2 \eTH \bTH Monolítico \eTH \eTR
\bTR \bTD 3 lotes, 7×7 \eTD \bTD 174 variables \eTD \bTD 1.191 variables, 1.725 restricciones \eTD \eTR
\bTR \bTD 6 lotes, 5×5 \eTD \bTD 369 variables \eTD \bTD 1.353 variables, 2.460 restricciones \eTD \eTR
\bTR \bTD 6 lotes, 7×7 \eTD \bTD 357 variables \eTD \bTD 2.391 variables, 3.486 restricciones \eTD \eTR
\bTR \bTD 6 lotes, 10×10 \eTD \bTD 357 variables \eTD \bTD 4.695 variables, 6.474 restricciones \eTD \eTR
\eTABLE
\blank[medium]

El split resuelve además una Etapa 1 por lote, que no está en la tabla.

{\bf Control de equivalencia.} En los 40 puntos de los estudios 08 y 10 donde los dos modelos
probaron el óptimo, 39 dieron exactamente el mismo objetivo. En el restante (3 lotes, 6×6) el monolítico
reportó 26.709,36 contra 26.709,60 del split, un 0,0009\% menos, con los dos declarando
optimalidad. La hipótesis, sin verificar, es numérica: con $M = 10^6$, un binario apenas
fraccionario dentro de la tolerancia de integralidad afloja una restricción big-M en cerca de una
unidad. Es la misma clase de problema de precisión que ya tienen los tests de OPP\_v2 y OPP\_v3.
\stopsection

\startsection[title={Estudio 08: la frontera del monolítico}]
\placefigure[force][fig:frontera]{Resultado del monolítico en cada punto de la frontera. Azul: probó el óptimo (tiempo total); rojo: cortado a los 10 minutos (brecha al corte).}
{%%FIG:frontera%%}

\blank[medium]
\bTABLE[split=yes, align=middle, frame=on, width=fit]
\setupTABLE[row][1][background=color, backgroundcolor=darkgray, style=bold]
\bTR \bTH Lotes \eTH \bTH Mayor grilla resuelta \eTH \bTH Primera grilla sin cerrar \eTH \bTH Split en ese punto \eTH \eTR
\bTR \bTD 1 y 2 \eTD \bTD 7×7 \eTD \bTD -- \eTD \bTD -- \eTD \eTR
\bTR \bTD 3 \eTD \bTD 6×6 (139 s) \eTD \bTD 7×7, brecha 11,4\% \eTD \bTD óptimo en 9,7 s \eTD \eTR
\bTR \bTD 4 \eTD \bTD 5×5 (47 s) \eTD \bTD 6×6, brecha 13,2\% \eTD \bTD óptimo en 8,4 s \eTD \eTR
\bTR \bTD 6 \eTD \bTD 3×3 (192 s) \eTD \bTD 5×5, brecha 3,0\% \eTD \bTD óptimo en 3,6 min \eTD \eTR
\eTABLE
\blank[medium]

Las celdas \quote{infactible} son instancias sin solución en los dos modelos: en 4×4 un lote de
26 órdenes pasa por casi todas las ubicaciones, y la restricción de un arco por vértice impide
el camino (ver el documento del split).

\placefigure[force][fig:ratio]{Cuántas veces más lento es el monolítico que el split sobre la misma instancia (escala logarítmica). Los círculos abiertos son cotas inferiores: el monolítico no terminó.}
{%%FIG:ratio%%}

\startitemize
\item {\bf Con un lote, el monolítico es igual o más rápido.} No hay nada que coordinar entre
lotes, y el split paga la preparación de dos etapas.
\item {\bf Con dos o más lotes, la desventaja crece con la grilla.} Con 2 lotes, desde 5×5,
el monolítico tarda entre 11 y 15 veces más; con 3 y 4 lotes llega a 27 veces antes de dejar de
cerrar, y en los puntos cortados es al menos 59 y 68 veces.
\item {\bf Con 6 lotes en una grilla chica, los dos modelos empatan} (177 s y 192 s en 3×3):
casi todo el costo es secuenciar, y eso lo pagan los dos. Pero con la misma cantidad de lotes en
5×5 el split sigue cerrando (3,6 min) y el monolítico no.
\stopitemize
\stopsection

\startsection[title={Encontrar no es el problema; demostrar sí}]
\placefigure[force][fig:trayectoria]{Evolución de la mejor solución del monolítico, como porcentaje sobre el óptimo, en los tres puntos donde no cerró.}
{%%FIG:trayectoria%%}

\blank[medium]
\bTABLE[split=yes, align=middle, frame=on, width=fit]
\setupTABLE[row][1][background=color, backgroundcolor=darkgray, style=bold]
\bTR \bTH Punto \eTH \bTH A menos del 1\% \eTH \bTH Llegó al óptimo \eTH \bTH Cota dual al corte \eTH \bTH Nodos \eTH \eTR
\bTR \bTD 3 lotes, 7×7 \eTD \bTD 41 s \eTD \bTD no \eTD \bTD 90,1\% del óptimo \eTD \bTD 45.924 \eTD \eTR
\bTR \bTD 4 lotes, 6×6 \eTD \bTD 29 s \eTD \bTD no \eTD \bTD 88,5\% \eTD \bTD 90.695 \eTD \eTR
\bTR \bTD 6 lotes, 5×5 \eTD \bTD 75 s \eTD \bTD 85 s \eTD \bTD 97,1\% \eTD \bTD 102.612 \eTD \eTR
\eTABLE
\blank[medium]

En los tres casos el monolítico tenía una solución a menos del 1\% del óptimo antes de los
75 segundos; con 6 lotes en 5×5 llegó al óptimo exacto a los 85 s. El resto de los 10 minutos
lo pasó sin poder demostrarlo: su cota dual quedó entre el 88\% y el 97\% del óptimo. En esas
mismas instancias, la Etapa 2 del split cerró con 27, 165 y 359.557 nodos.

Tampoco es que la relajación del monolítico arranque peor. En casi todas las instancias, su
cota en el nodo raíz es parecida a la de la Etapa 2 del split:

\placefigure[force][fig:raiz]{Cota dual en el nodo raíz, como porcentaje del óptimo: Etapa 2 del split contra monolítico, sobre la misma instancia.}
{%%FIG:raiz%%}

La debilidad de la relajación es común: viene del secuenciamiento con big-M, que los dos modelos
comparten. Lo que cambia es el costo de cerrarla. El monolítico tiene entre 2 y 13 veces más
variables que la Etapa 2 del split, cada nodo le cuesta más, y además de secuenciar tiene que ramificar sobre las rutas.
Hay tres excepciones claras, donde su cota raíz es mucho peor: la grilla 10×10 (39\% contra 81\%),
las velocidades en bandas (41\%–53\% contra 76\%–87\%) y 6 lotes en 5×5 (48\% contra 87\%). En
esos casos pesa justo el término que distingue a los dos modelos: en el monolítico, el tiempo de
ejecución (parte fija más ruta dividida por la velocidad del picker) es una variable atada con
big-M; en el split es un dato. Con grillas grandes la ruta es larga, con velocidades distintas
el término cambia por picker, y con muchos lotes hay más pares picker-lote. Es una explicación
consistente con los datos, pero no está verificada.
\stopsection

\startsection[title={Estudio 09: las mismas instancias que el split}]
Las 15 instancias de los estudios 01, 02 y 03 del split, con el mismo límite de 300 s. El split
obtuvo exactamente los mismos óptimos que en esos estudios, lo que confirma que son las mismas
instancias.

\placefigure[force][fig:compartidos]{Tiempo del split y del monolítico en los casos compartidos (escala logarítmica). o: órdenes por lote; p: pickers; g10: grilla 10×10; w: pesos 2/6 con esas órdenes por lote; vd: velocidades distintas; b2 y b3: dos y tres bandas. Sobre la barra del monolítico, su brecha al corte, o s/sol. si no encontró ninguna solución.}
{%%FIG:compartidos%%}

\startitemize
\item {\bf El split resolvió las 15, entre 13 y 167 s. El monolítico no demostró ninguna.}
\item {\bf En 7 de las 15, el monolítico no encontró ni una solución factible en 5 minutos:}
22 y 26 órdenes por lote, la grilla 10×10, pesos 2/6 con 18 órdenes, velocidades distintas y
las dos configuraciones en bandas. Con 6 lotes en 7×7, empieza a costarle también encontrar.
\item {\bf Cuando encuentra, encuentra bien:} en los otros 8 casos su mejor solución quedó entre
0,3\% y 1,9\% sobre el óptimo (4,9\% en uno).
\item {\bf Las brechas grandes son de la cota, no de la solución:} con 3 a 5 pickers la cota
dual queda entre el 49\% y el 59\% del óptimo, y la brecha reportada llega al 105\% aunque la solución
esté a menos del 1\%.
\stopitemize

\placefigure[force][fig:cotas]{Dónde terminan las dos cotas del monolítico al cabo de 300 s, como proporción del óptimo del split.}
{%%FIG:cotas%%}
\stopsection

\startsection[title={Estudio 10: qué factor le pega más al monolítico}]
Las preguntas de los estudios 01, 02, 03 y 07 del split, a una escala que el monolítico
resuelve: base de 4 lotes, 2 pickers, 26 órdenes por lote, grilla 5×5, pesos 4/4 y velocidades
uniformes, moviendo un factor por vez. Tope de 10 minutos por modelo. En la base, el split
tarda 7 s y el monolítico 17 s (2,4 veces). Corrió en paralelo con la suite de stress del split
primero y con la del monolítico después.

\placefigure[force][fig:m10_grilla]{Tiempo de los dos modelos según la grilla (escala logarítmica).}
{%%FIG:m10_grilla%%}

\placefigure[force][fig:m10_velocidades]{Tiempo de los dos modelos según las velocidades de la cuadrilla (escala logarítmica). 2P, 3P y 4P: cantidad de pickers.}
{%%FIG:m10_velocidades%%}

\placefigure[force][fig:m10_ordenes]{Tiempo de los dos modelos según las órdenes por lote, con pesos 4/4 y 2/6 (escala logarítmica).}
{%%FIG:m10_ordenes%%}

\placefigure[force][fig:m10_lotes]{Tiempo de los dos modelos según la cantidad de lotes (escala logarítmica).}
{%%FIG:m10_lotes%%}

\placefigure[force][fig:m10_pickers]{Tiempo de los dos modelos según la cantidad de pickers (escala logarítmica).}
{%%FIG:m10_pickers%%}

\placefigure[force][fig:m10_le1]{Tiempo de los dos modelos con y sin el límite de un arco por vértice (escala logarítmica).}
{%%FIG:m10_le1%%}

\blank[medium]
\bTABLE[split=yes, align=middle, frame=on, width=fit]
\setupTABLE[row][1][background=color, backgroundcolor=darkgray, style=bold]
\bTR \bTH Factor \eTH \bTH Rango \eTH \bTH Monolítico / split \eTH \eTR
\bTR \bTD Grilla \eTD \bTD 3×3 a 7×7 \eTD \bTD de 1,1 a 33 veces; en 7×7 no cierra (brecha 15,3\%) \eTD \eTR
\bTR \bTD Velocidades \eTD \bTD 3 pickers, 2 y 3 bandas \eTD \bTD 52 veces; con 3 bandas no cierra (brecha 1,7\%) \eTD \eTR
\bTR \bTD Órdenes por lote \eTD \bTD 10 a 50 \eTD \bTD 188 veces con 10, 64 con 14, entre 4 y 10 desde 18 \eTD \eTR
\bTR \bTD Lotes \eTD \bTD 2 a 6 \eTD \bTD de 1,2 a 8 veces \eTD \eTR
\bTR \bTD Pickers \eTD \bTD 1 a 4 \eTD \bTD entre 2,4 y 6,5 veces \eTD \eTR
\bTR \bTD Pesos 2/6 \eTD \bTD 10, 26 y 50 órdenes \eTD \bTD sin efecto sistemático sobre el monolítico \eTD \eTR
\bTR \bTD Sin $\leq 1$ \eTD \bTD 5×5 y 6×6 \eTD \bTD el monolítico tarda el doble en 5×5 y la mitad en 6×6 \eTD \eTR
\eTABLE
\blank[medium]

\startitemize
\item {\bf La grilla es el factor principal.} En 3×3 los dos modelos empatan; en 6×6 el
monolítico ya es 33 veces más lento y en 7×7 no cierra, mientras el split tarda 14 s. Es el
mismo eje que limita a la Etapa 1 del split, pero el monolítico lo sufre mucho antes.
\item {\bf Las velocidades en bandas golpean al monolítico, no al split, a esta escala.} Con 3
pickers, el split tarda lo mismo con o sin bandas (3 a 4 s); el monolítico pasa de 18 s a 146 s
con 2 bandas y no cierra con 3. Con 4 pickers para 4 lotes (un lote por picker, sin
secuenciamiento) el efecto casi desaparece. Con velocidades todas distintas y 2 pickers, apenas
2,9 veces.
\item {\bf Pocas órdenes por lote lo hunden.} Con 10 órdenes por lote el monolítico tarda 413 s
contra 2 s del split. Es el efecto de objetivo aplanado que en el split encarecía la Etapa 2,
mucho más fuerte en el monolítico.
\item {\bf Lotes y pickers pesan menos}, y los pesos 2/6 no cambian nada de forma sistemática.
\item {\bf Quitar el límite de un arco por vértice no es una palanca,} igual que en el split:
en 5×5 el monolítico tarda el doble, en 6×6 la mitad, siempre con el mismo óptimo que el split.
\stopitemize
\stopsection

\startsection[title={Tamaño y memoria no son el límite}]
\placefigure[force][fig:tamano]{Tiempo del monolítico contra la cantidad de variables del modelo (escala logarítmica en los dos ejes).}
{%%FIG:tamano%%}

Los modelos son chicos: el mayor cortado tiene 4.695 variables. Con 10 minutos la memoria pico
no pasó de 302 MB; con 40 minutos (estudio 11) llegó a 636 MB, porque el árbol de búsqueda sigue
creciendo. La compilación con ZIMPL tardó menos de un segundo en todos los puntos. Hay modelos de
1.191 variables que no cierran y de 2.199 que encuentran buenas soluciones: lo que decide es la
estructura, no el tamaño. La única excepción posible son los escenarios heredados de 100×100
de la suite de stress, donde el monolítico terminó sin resultado o con error; su memoria no se
midió.
\stopsection

\startsection[title={Estudio 11: con 40 minutos}]
Los tres puntos donde el monolítico dejó de cerrar a los 10 minutos, más la instancia exacta del
test de stress 04, con el tope de la suite: 40 minutos.

\blank[medium]
\bTABLE[split=yes, align=middle, frame=on, width=fit]
\setupTABLE[row][1][background=color, backgroundcolor=darkgray, style=bold]
\bTR \bTH Punto \eTH \bTH Split \eTH \bTH Monolítico, 10 min \eTH \bTH Monolítico, 40 min \eTH \eTR
\bTR \bTD 3 lotes, 7×7 \eTD \bTD óptimo en 8 s \eTD \bTD brecha 11,4\% \eTD \bTD brecha 11,0\% \eTD \eTR
\bTR \bTD 4 lotes, 6×6 \eTD \bTD óptimo en 10 s \eTD \bTD brecha 13,2\% \eTD \bTD brecha 10,4\% \eTD \eTR
\bTR \bTD 6 lotes, 5×5 \eTD \bTD óptimo en 3,9 min \eTD \bTD brecha 3,0\% \eTD \bTD óptimo en 32 min \eTD \eTR
\bTR \bTD test 04 (6 lotes, 7×7) \eTD \bTD óptimo en 3,0 min \eTD \bTD -- \eTD \bTD brecha 143,5\% \eTD \eTR
\eTABLE
\blank[medium]

\placefigure[force][fig:m11_cotas]{Cota dual del monolítico como porcentaje del óptimo: en el nodo raíz, a los 10 minutos y a los 40.}
{%%FIG:m11_cotas%%}

\placefigure[force][fig:m11_trayectoria]{Evolución de la mejor solución del monolítico durante los 40 minutos, como porcentaje sobre el óptimo (escala logarítmica en el tiempo).}
{%%FIG:m11_trayectoria%%}

\startitemize
\item {\bf Cuadruplicar el tiempo casi no mueve la cota.} En 3 lotes en 7×7 y 4 lotes en 6×6
quedó en el 90\% del óptimo, igual que a los 10 minutos, aunque el monolítico recorrió seis y
cuatro veces más nodos. Las brechas bajaron poco: de 11,4\% a 11,0\% y de 13,2\% a 10,4\%.
\item {\bf Las soluciones ya estaban.} En 3 lotes en 7×7 la mejor solución quedó a 0,05\% del
óptimo; en 4 lotes en 6×6 encontró el óptimo exacto a los 972 s y no pudo demostrarlo en los 24
minutos que le quedaban.
\item {\bf Solo 6 lotes en 5×5 cerró}: el monolítico demostró el óptimo a los 32 minutos, 8 veces
lo que tardó el split. Ahí la cota raíz era la más débil (48\%), pero la búsqueda sí la subió.
\item {\bf Test 04: una solución casi óptima y una brecha del 143,5\%.} A los 511 s ya tenía una
solución a menos del 1\% del óptimo, y terminó a 0,66\%. Pero su cota dual apenas pasó del 39\%
en la raíz al 41\% a los 40 minutos. Es el mismo resultado que M0 a los 25 minutos (144,3\%): con
esta instancia, el tiempo extra no cambia nada.
\stopitemize
\stopsection

\startsection[title={Predicción y resultados de la suite}]
Esta predicción se escribió con la suite de stress del monolítico corriendo y sin conocer sus
resultados.

\startitemize
\item {\bf Test 04 (6 lotes, 7×7, perfil realista):} es la misma instancia que M0. En 25
minutos, con otros tres puntos corriendo en paralelo, M0 terminó con una solución a menos del
1\% del óptimo (79.822,91 contra 79.046,86) y una brecha del 144\%. Predicción: {\bf no termina en 40 minutos}, probablemente con una buena
solución y una brecha grande.
\item {\bf Tests 05 a 19:} todos tienen 6 o más lotes en grillas de 7×7 o mayores, o cuadrillas
más grandes. El monolítico ya no cierra 6 lotes en 5×5 y, en la mitad de los casos con 6 lotes en
7×7, no encuentra solución en 5 minutos. Predicción: {\bf ninguno termina}, y en varios puede no
haber ni una solución que reportar.
\item {\bf Tests heredados 01 a 03:} quedaron fuera de la predicción. Del 01 (100×100), la
corrida anterior del monolítico terminó sin resultado.
\stopitemize

{\bf Lo que confirmó el estudio 11:} sobre la instancia exacta del test 04, el monolítico llega a
los 40 minutos con una solución a 0,66\% del óptimo y una brecha del 143,5\%. La suite no registra
brechas, solo el timeout, así que ese número es la predicción cuantitativa para el test 04.

{\bf Resultados de la suite:} el monolítico terminó %%NUM:stress_mono_ok%% de los
%%NUM:stress_total%% tests en 40 minutos; el split resolvió %%NUM:stress_split_ok%%, ninguno en
más de %%NUM:stress_split_max%% minutos. Los tiempos del split son la duración completa del test
(compilación, Etapa 1 con seis procesos y Etapa 2) y son un techo: la máquina entró en suspensión
durante parte de esa corrida, y el tiempo suspendido cuenta en la duración pero no en el tope.

%%TABLE:stress%%

Los tests que el split tampoco termina son los del perfil realista con 8 lotes o más (12 a 14 y
16 a 19), la grilla de 16×16 (08) y dos de los heredados (01 y 03). En esos casos los dos modelos llegan al
tope; en todos los demás, la diferencia es la que ya mostraban los estudios: el split cierra en
minutos instancias donde el monolítico no demuestra el óptimo en 40.
\stopsection

\startsection[title={Síntesis}]
\startitemize[n]
\item {\bf El monolítico deja de cerrar mucho antes que el split:} con el perfil realista, a
los 10 minutos llega hasta 3 lotes en 6×6, 4 lotes en 5×5 y 6 lotes en 3×3. El split cierra
todos esos puntos en segundos o pocos minutos.
\item {\bf Su desventaja crece con la grilla y con los lotes:} de igual a 15 veces más lento con
1 y 2 lotes; con 3 y 4, hasta 27 veces cuando todavía cierra y al menos 59 a 68 cuando no.
\item {\bf Encuentra buenas soluciones rápido y no las puede demostrar:} a menos del 1\% del
óptimo en segundos, con la cota dual estancada entre el 39\% y el 97\%.
\item {\bf La cota raíz es parecida a la del split}, salvo con velocidades en bandas, grillas
grandes o muchos lotes. La diferencia está en cuánto le cuesta cerrarla: más variables por nodo
y rutas que el split ya fijó.
\item {\bf Lo que más lo castiga es la grilla, las velocidades en bandas y las pocas órdenes por
lote.} Lotes y pickers pesan menos, y ni los pesos asimétricos ni quitar el límite de un arco
por vértice lo ayudan de forma sistemática.
\item {\bf Más tiempo casi no ayuda:} de 10 a 40 minutos la cota dual no se movió en dos de los
tres puntos que no cerraban, y en el test 04 pasó del 39\% al 41\%. La brecha de una corrida corta
predice bien la de una larga.
\item {\bf No es un problema de tamaño ni de memoria:} los modelos tienen pocos miles de
variables y usan menos de 650 MB aun a los 40 minutos.
\item {\bf En la suite de stress:} el monolítico termina %%NUM:stress_mono_ok%% de
%%NUM:stress_total%% tests en 40 minutos, contra %%NUM:stress_split_ok%% del split.
\stopitemize
\stopsection

\stoptext
